\subsection{Finite Refinement}
\label{app:refinement-convergence}

All arguments in this section concern a fixed finite KL-UOT problem in exact arithmetic.
Reference masses and penalties are positive, costs are finite and nonnegative, and each active support covers every node.
The hierarchy has finitely many levels, so finite termination at each level implies finite termination of the hierarchy.
Throughout, pricing is exact bidirectional top-$\gamma$ with $\gamma\geq1$, sparsification preserves transported marginals and cost, and pruning retains positive-flow and newly priced edges.
A restricted KKT point extended by zero is globally optimal whenever its potentials satisfy every full-space dual inequality, since these supply exactly the missing KKT conditions.

\subsubsection{Finite Termination with Inexact Solves}
\label{app:refinement-inexact}

For an edge $e=(i,j)$, define its gain at transported marginals $(r,q)$ by
\begin{equation*}
g_e(r,q)=\max_{t\geq0}\left\{-C_{ij}t-\rho_a[\KL(r+t\mid a_i)-\KL(r\mid a_i)]-\rho_b[\KL(q+t\mid b_j)-\KL(q\mid b_j)]\right\}.
\end{equation*}
This is a feasible primal improvement even when $r=0$ or $q=0$.
For positive marginals, the primal-consistent score is $\sigma^\mathbf X_{ij}=-C_{ij}-\rho_a\log(r_i/a_i)-\rho_b\log(q_j/b_j)$.
At a zero marginal, the corresponding score is interpreted as $+\infty$ for the purpose of the one-sided coordinate direction.

Let $\mathbf X^k$ be the returned coupling, $\Delta_k=\Phi(\mathbf X^k)-\Phi(\mathbf X^{k+1})$, and define the feasible coordinate gain
\begin{equation*}
G_k=\max_{e=(i,j)\in\mathcal V_k\setminus\mathcal N_k}g_e(r_i^k,q_j^k),
\end{equation*}
with $G_k=0$ for an empty candidate set.
Write $v_k=\max_{i,j}[\sigma^k_{ij}]_+$ and $a_k=\max_{(i,j)\in\mathcal N_k}[\sigma^k_{ij}]_+$ for full and active score maxima.
The one-sided marginal-gradient error is
\begin{equation}
\label{eq:refinement-fidelity}
\begin{aligned}
\Omega_k={}&\max_i[f_i^k+\rho_a\log(r_i^k/a_i)]_+\\
&+\max_j[g_j^k+\rho_b\log(q_j^k/b_j)]_+,
\end{aligned}
\end{equation}
where a zero transported marginal contributes zero to its one-sided error.

\begin{theorem}[Inexact refinement with explicit output conditions]
\label{thm:refinement-explicit}
Use $\ell_2$ stopping measures with a restricted tolerance $\varepsilon_{\mathrm r}>0$ and a full-space tolerance $\tau>0$.
Set $d_k=\varepsilon_{\mathrm r}^2(1+|\Phi(\mathbf X^k)|)$.
Suppose each restricted call returns in finite time with all three restricted stopping measures at most $\varepsilon_{\mathrm r}$, and consecutive refinement outputs satisfy
\begin{equation}
\label{eq:refinement-output-conditions}
\begin{gathered}
\Delta_k\geq G_k/2,\\
\Delta_k\geq d_k
\quad\text{or}\quad
\max\{a_{k+1},\Omega_{k+1}\}\leq v_k/2.
\end{gathered}
\end{equation}
Then refinement terminates after finitely many rounds with restricted stopping measures at most $\varepsilon_{\mathrm r}$ and full-space violation $\eta\leq\tau$.
In particular, $\varepsilon_{\mathrm r}=\tau=\varepsilon$ gives the KKT stopping tolerance used in Algorithm~\ref{alg:ours}.
\end{theorem}

\begin{lemma}[Positive gain at a fixed score margin]
\label{lem:refinement-uniform-gain}
Fix a finite objective sublevel set and a margin $\delta>0$.
There is a constant $\Gamma_\delta>0$ such that every edge with primal-consistent score at least $\delta$ has coordinate gain at least $\Gamma_\delta$.
\end{lemma}

\begin{proof}
Nonnegative transport costs and coercivity of scalar KL bound each transported marginal on the sublevel set.
For a fixed edge, the score-margin condition is equivalently
\begin{equation*}
r^{\rho_a}q^{\rho_b}\leq a_i^{\rho_a}b_j^{\rho_b}\exp(-C_{ij}-\delta),
\end{equation*}
which defines a closed set, including zero marginals, within these bounded marginal intervals.
The coordinate minimizer lies in a common finite interval: for $t>\exp((\rho_a\log a_i+\rho_b\log b_j-C_{ij})/(\rho_a+\rho_b))$, the derivative of the coordinate objective is positive for every $r,q\geq0$.
Continuity of scalar KL on nonnegative arguments therefore makes $g_e(r,q)$ continuous on the compact marginal domain.
On the score-margin set it is strictly positive, since the initial derivative is at most $-\delta$, or is negative infinite at a zero marginal.
Its minimum on every nonempty edgewise score-margin set is thus positive.
Taking the minimum over the finitely many such edgewise sets proves the claim.
\end{proof}

\begin{proof}[Proof of Theorems~\ref{thm:refinement-explicit} and~\ref{thm:refinement}]
Extend all couplings by zero to the full edge set.
The first condition in~\eqref{eq:refinement-output-conditions} and $G_k\geq0$ give objective nonincrease, and sparsification and pruning preserve the current objective.
All subsequent couplings therefore lie in the sublevel set determined by the first returned coupling $\mathbf X^0$.
Whenever the sufficient-decrease alternative holds, the objective falls by at least $\varepsilon_{\mathrm r}^2$.
There can be only finitely many such rounds because $\Phi\geq0$.

Suppose, for contradiction, that the refinement produces an infinite nonterminal sequence.
After finitely many rounds, every output must satisfy the pricing-fidelity alternative.
Nontermination under~\eqref{eq:hierarchy-full-violation}, with $p=2$, implies
\begin{equation*}
v_k>b_\tau,\qquad b_\tau=\frac{\tau(1+\|\mathbf C\|_F)}{\sqrt{mn}}>0.
\end{equation*}
Consider a fidelity-controlled output at round $k+1$.
If $v_{k+1}<3v_k/4$, the largest raw score contracts geometrically on this round.
Otherwise, $a_{k+1}\leq v_k/2<v_{k+1}$ ensures that any full-space maximizer is inactive, and exact nodewise top-$\gamma$ pricing retains at least one such edge $e$.
For its endpoints with positive transported marginals,~\eqref{eq:refinement-fidelity} gives
\begin{equation*}
\sigma_e^{\mathbf X^{k+1}}\geq\sigma_e^{k+1}-\Omega_{k+1}\geq v_k/4>b_\tau/4.
\end{equation*}
If either transported marginal is zero, the same lower bound holds with the extended one-sided score.
Lemma~\ref{lem:refinement-uniform-gain} supplies $G_{k+1}\geq\Gamma_{b_\tau/4}>0$.
Protection of newly priced edges and the current positive support preserves the feasible coordinate improvement after pruning.
The following restricted solve consequently decreases the objective by at least $\Gamma_{b_\tau/4}/2$.
There can be only finitely many such noncontracting rounds.

After the last sufficient-decrease or noncontracting round, the raw maxima would therefore satisfy $v_{k+1}<3v_k/4$ forever.
Since $\eta_k\leq\sqrt{mn}\,v_k/(1+\|\mathbf C\|_F)$, this forces $\eta_k\leq\tau$ after finitely many more rounds, a contradiction.
Each restricted call returns in finite time by hypothesis and satisfies the prescribed restricted stopping measures, proving the stated finite-tolerance termination.
\end{proof}

\paragraph{Exact restricted solves.}
If every call returns an exact restricted primal--dual optimum, a positive dual score gives a one-edge strict descent direction.
Sparsification preserves the objective and marginals, and retaining that edge and the positive support makes the trial feasible on the next support.
Thus the restricted optimal value strictly decreases at each nonterminal round and cannot revisit any of the finitely many supports.
With exact full-space feasibility stopping, refinement therefore terminates at a global primal--dual optimum.
This idealized result assumes an exact oracle; the numerical algorithm uses the finite-tolerance result above.

\subsubsection{Attainability of the Stopping Conditions}
\label{app:refinement-stop-attainability}

The preceding theorem assumes finite restricted returns satisfying~\eqref{eq:refinement-output-conditions}.
The following result connects that assumption to \subsolver and the stopping rule in Appendix~\ref{app:refinement-local-stopping}.

\begin{proposition}[Finite attainability and preservation]
\label{prop:refinement-stop-attainability}
Consider a nonterminal same-level support update, with padded anchor $\mathbf x^0$, previous violation maximum $v>0$, and candidate gain $G$.
Suppose the updated support retains the anchor's positive support and all candidate edges, and \subsolver satisfies the assumptions of Theorem~\ref{thm:pduot-global}.
At any fixed restricted tolerance $\varepsilon_{\mathrm r}>0$, the complete stopping rule, including the optimal-anchor readout, accepts after finitely many PDHG evaluations and implies~\eqref{eq:refinement-output-conditions}.
Subsequent sparsification preserving marginals and cost preserves the restricted certificates on the original active graph after refreshing bound multipliers, and preserves these output conditions.
Pruning zero-flow edges while protecting newly priced edges preserves the current coupling and the candidate coordinate trials.
\end{proposition}

\begin{proof}
\textit{Decrease guarantee.}
Using the notation of Appendix~\ref{app:refinement-local-stopping}, the generalized-KL identity and projection complementarity give
\begin{equation*}
\Phi(\mathbf x^0)-\Phi(\mathbf x)
=L+\sum_{e\in\mathcal N}[-w_e(x_e-x_e^0)]_+\geq L.
\end{equation*}
The implemented tests~\eqref{eq:refinement-implemented-gate} therefore imply~\eqref{eq:refinement-output-conditions}.

\par\smallskip\noindent\textit{Finite attainability.}
Theorem~\ref{thm:pduot-global} gives convergence of mapped iterates and projection multipliers to a finite restricted KKT point with coupling $\mathbf x^\star$.
Its marginals are positive by node coverage, so all three restricted stopping measures and $a,\Omega$ tend to zero.
Moreover, $\mathbf w\to0$ and the bound $L$ tends to $\Delta^\star:=\Phi(\mathbf x^0)-\Phi(\mathbf x^\star)$.
Every candidate coordinate trial is feasible on the updated support, giving $\Delta^\star\geq G$.
If $G>0$, then $L\geq G/2$ eventually holds with strict margin.
If $G=0$ and the anchor is suboptimal, $\Delta^\star>0$ likewise gives $L\geq0$ eventually.
In either case, $\max\{a,\Omega\}\leq v/2$ and the positive-tolerance restricted tests eventually hold, so the call accepts after finitely many evaluations.
If the anchor is already optimal, node coverage gives positive anchor marginals and its marginal-induced potentials satisfy restricted KKT conditions.
The initial optimal-anchor readout therefore accepts with $G=0$, zero restricted residuals, and $a=\Omega=0$.
This handles zero limiting decrease without requiring the iterative bound $L$ to cross zero.

\par\smallskip\noindent\textit{Preservation under support operations.}
Sparsification leaves marginals, transport cost, potentials, and the active graph unchanged while redistributing flow onto a subset of the previous positive edges.
Thus it preserves the objective decrease, $\Omega$, both positive-score maxima, the marginal residual, and the restricted split gap.
Retain zero multipliers on positive edges and choose the closest nonnegative multiplier to the dual slack on each zero edge.
This cannot increase the stationarity residual: previously positive edges had zero multipliers, and previously zero edges already had admissible nonnegative multipliers.
Consequently the restricted tests and~\eqref{eq:refinement-output-conditions} remain satisfied.
Later pruning removes only zero-flow edges, preserving the coupling, marginals, objective, and full-space potentials.
The accepted restricted certificates refer to the original active graph, as specified in Appendix~\ref{app:hierarchy-pruning}; subsequent restricted solves check their own updated supports.
Keeping the positive support and newly priced edges also preserves all coordinate trials used by the finite-refinement proof.
\end{proof}
